% part: intuitionistic-logic
% chapter: introduction
% section: constructive-reasoning

\documentclass[../../../include/open-logic-chapter]{subfiles}

\begin{document}

\olfileid{int}{int}{cr}

\section{构造性推理}

与通过模态算子或二阶量词来扩展经典逻辑不同，直觉主义逻辑的“非经典”性体现在它\emph{限制}了经典逻辑。经典逻辑在许多方面是\emph{非构造性的}。直觉主义逻辑旨在捕捉一种更具“构造性”的推理方式，这种推理方式是构造性数学的特征。下面的例子或许可以说明其背后的一些动机。

假设有人声称，他们找到了一个具有如下性质的自然数~$n$：如果 $n$ 是偶数，那么 黎曼猜想为真；如果 $n$ 是奇数，那么 黎曼 猜想为假。好消息！{}黎曼 猜想是否为真是数学的重大悬而未决的问题之一，而他们似乎已经把这个问题归结为了一个计算问题，即，判定一个特定数是偶数还是奇数。

这个神奇的~$n$ 是多少呢？他们是这样描述的：$n$ 是这样一个自然数：若 黎曼 猜想为真，则它等于 $2$，否则等于 $3$。

你愤怒地要求退款。从经典的角度看，上述描述确实唯一地确定了~$n$ 的一个值；但你真正想要的是一个能\emph{显式地}给出的~$n$ 的值。

再举一个也许不那么牵强的例子，考虑下面的问题。我们知道，一个无理数取有理数次方，有可能得到有理数的结果。例如，$\sqrt{2}^2 = 2$。不那么清楚的是，一个无理数取\emph{无理}数次方，是否可能得到有理数的结果。下面的定理给出了肯定的回答：

\begin{thm}
存在无理数 $a$ 和 $b$，使得 $a^b$ 是有理数。
\end{thm}

\begin{proof}
考虑 $\sqrt{2}^{\sqrt{2}}$。如果它是有理数，我们就完成了证明：取 $a = b = \sqrt{2}$ 即可。否则，它是无理数。那么我们有
\[
(\sqrt{2}^{\sqrt{2}})^{\sqrt{2}} = \sqrt{2}^{\sqrt{2} \cdot
  \sqrt{2}} = \sqrt{2}^2 = 2,
\]
而这是有理数。所以，在这种情况下，令 $a$ 为 $\sqrt{2}^{\sqrt{2}}$，令 $b = \sqrt 2$。
\end{proof}

这算是一个有效的证明吗？大多数数学家认为是。但同样，这里还是有一点让人不满足：我们证明了存在一对具有某种性质的实数，却无法说出究竟是哪一对数。其实，我们可以用另一种方式证明同样的结果，使得证明中数对 $a$、$b$ \emph{是}被给出的：取 $a = \sqrt{3}$ 和 $b = \log_3 4$。那么
\[
a^b = \sqrt{3}^{\log_3 4} = 3^{1/2 \cdot \log_3 4} = (3^{\log_3
  4})^{1/2} = 4^{1/2}= 2,
\]
因为 $3^{\log_3 x} = x$。

直觉主义逻辑旨在捕捉一种排除了第一种证明中那种推理方式的推理。证明存在一个满足 $!A(x)$ 的~$x$，意味着你必须给出一个具体的~$x$，并证明它满足 $!A$，就像在第二个证明中所做的那样。证明 $!A$ 或 $!B$ 成立，则要求你能证明其中之一。

从形式上说，直觉主义逻辑就是以某种方式限制经典逻辑的推导系统所得到的结果。从数学的角度看，这些都只是形式演绎系统，但正如已经指出的，它们旨在捕捉一种数学推理方式。人们可以认为，这对应于某种数学哲学观点（例如 Brouwer（布劳威尔）的直觉主义）下合理的推理方式；也可以认为，这是一种更“具体”、更令人满意的数学推理（类似于 Bishop（毕晓普）的构造主义）；当然，人们也可以争论形式描述是否抓住了非形式的动机。但无论我们持何种哲学立场，我们都可以把直觉主义逻辑作为一个形式呈现的逻辑系统来研究；并且由于各种原因，许多数理逻辑学家觉得这样做很有趣。

\end{document}